% Standalone source. Compile twice with pdflatex or use latexmk -pdf.
% All text, tables, proof details, code and references are embedded.
\documentclass[11pt,a4paper]{article}
\usepackage[T1]{fontenc}
\usepackage[utf8]{inputenc}
\usepackage{lmodern,microtype,amsmath,amssymb,amsthm,mathtools}
\usepackage[left=24mm,right=24mm,top=23mm,bottom=24mm,headheight=15pt,headsep=6mm]{geometry}
\usepackage{booktabs,array,tabularx,xcolor,fancyhdr,enumitem}
\usepackage{hyperref,bookmark}
\definecolor{ink}{HTML}{183F48}
\definecolor{muted}{HTML}{42666C}
\hypersetup{colorlinks=true,linkcolor=ink,citecolor=ink,urlcolor=ink,
pdftitle={Transcendence criteria from congruences and continued fractions},
pdfsubject={Diophantine approximation, Mahler values, and modular rotation codings}}
\newcommand{\Q}{\mathbb Q}\newcommand{\Z}{\mathbb Z}\newcommand{\PP}{\mathbb P}\newcommand{\SL}{\mathrm{SL}}
\newtheorem{theorem}{Theorem}[section]
\newtheorem{lemma}[theorem]{Lemma}
\newtheorem{proposition}[theorem]{Proposition}
\newtheorem{corollary}[theorem]{Corollary}
\setcounter{tocdepth}{1}
\setlist{nosep,topsep=4pt}
\setlength{\emergencystretch}{1.5em}
\setlength{\parindent}{1.1em}
\raggedbottom
\pagestyle{fancy}\fancyhf{}
\fancyhead[L]{\small\sffamily\color{muted} Congruences and continued fractions}
\fancyhead[R]{\small\sffamily\color{muted}\thepage}
\renewcommand{\headrulewidth}{0.25pt}
\begin{document}
\thispagestyle{empty}
\noindent{\small\sffamily\color{muted} COMPANION TO THE ONE-HUNDRED-CASE CATALOGUE\hfill 15 November 2026}\par
\vspace{1.25em}
\noindent{\fontsize{23}{27}\selectfont\bfseries\color{ink} Transcendence criteria from\\[1pt] congruences and continued fractions\par}
\vspace{.6em}
\noindent{\large Quadratic norms, Mahler values and modular rotation codings\par}
\vspace{.7em}
\noindent\rule{\linewidth}{.45pt}
\smallskip

\noindent\textbf{Abstract.}
A large power of seven dividing $p^2-2q^2$ can turn a rational approximation
at the critical real exponent into a proof of transcendence. We give a
bounded-partial-quotient construction and apply the criterion to ten explicit
lacunary quotients built from conjugate arguments in $\Q(\sqrt2)$. Mahler's method establishes algebraic
independence of all ten quotients and of their twenty underlying values.
A quantitative criterion using two adjacent continued-fraction convergents
then detects a computable number of irrationality exponent two even when
every fixed one-fraction test considered here fails to achieve a positive
power saving. For Wolf's prime continued fraction we prove the required
four-vector independence and identify the remaining congruence estimate.
For rotation codings built from Ramanujan's tau function we prove
transcendence for all integer bases, moduli and alphabet sizes at least
two. We classify the analogous codings for five Eisenstein series and
give exact computations at the finite places used in the proofs.
\smallskip

\noindent\textbf{Relation to the earlier paper.}
The earlier catalogue \cite{catalogue} starts from a transcendental parameter
$T$ and recovers a target as $R(T)$ for a nonconstant rational function $R$
with algebraic coefficients. The present problem is to obtain the arithmetic
input when a target is specified by a sequence. The proofs use divisibility of a norm of a rational approximant,
the number of independent convergent pairs at one prime, and the factor
complexity of a sequence of rotation digits. The catalogue identifiers in
Appendix~\ref{sec:rankings} locate the review that motivated these investigations.

\begingroup\small
\tableofcontents
\endgroup
\clearpage


\section{A split quadratic norm and one rational approximation}\label{sec:pell}
For a prime $p$ normalize $v_p(p)=1$. A split binary quadratic form can
record the quality of approximation at a finite place in an ordinary
integer. We use $P(s,t)=s^2-2t^2$ at $7$. Its two simple roots modulo
$7$ are $3$ and $-3$.

\begin{theorem}\label{th:pelldrop}
Let $x$ be a real number. Suppose infinitely many distinct reduced fractions
$a/b\ne x$, with $b>0$ and $a/b\to x$, satisfy, for some fixed
$\delta>0$,
\[
 b^2|x-a/b|\,7^{-v_7(a^2-2b^2)}\le b^{-\delta}.
\]
Then $x$ is transcendental. In particular it suffices that
$|x-a/b|\le C b^{-2}$ and $7^{v_7(a^2-2b^2)}\ge b^\eta$
on the same infinite sequence, for fixed $C,\eta>0$.
\end{theorem}
\begin{proof}
If $v_7(a^2-2b^2)=0$ on an infinite subsequence, Roth's theorem
\cite{pellroth} excludes an algebraic irrational $x$ there. For rational
$x=r/s$ the elementary bound $|x-a/b|\ge1/(sb)$ suffices in this case.
Otherwise take an infinite subsequence on which $7\mid a^2-2b^2$.
Then $7\nmid b$. Write $\rho^2=2$ in $\mathbb Z_7$ with
$\rho\equiv3\pmod 7$ by Hensel's lemma.
Exactly one of $a-\rho b$ and $a+\rho b$ is divisible by $7$.
After selecting one sign, a fixed $\alpha\in\{\rho,-\rho\}$ satisfies
\[
 |a/b-\alpha|_7=7^{-v_7(a^2-2b^2)}.
\]
The height $\max(|a|,b)$ is comparable to $b$ since $a/b\to x$.
If $x$ were algebraic, Roth--Ridout at the two places $\infty$ and $7$,
with targets $x$ and $\alpha$, would bound the product of the two errors
below by $H(a/b)^{-2-\delta/2}$ for all sufficiently large heights
\cite[Thm.~1]{cor}. The displayed hypothesis contradicts this bound.
Absorbing the fixed $C$ gives the last assertion.
\end{proof}

\subsection{Bounded partial quotients at the real critical exponent}
The fraction $[0;2,2,1,2,1,2,2,1]=101/239$ has
$101^2-2\cdot239^2=-104041$ and $v_7(-104041)=1$.
Arbitrarily deep divisibility can also be prescribed while the
partial quotients stay in this same two-letter alphabet.
\begin{theorem}\label{th:pellbounded}
There exists a computable transcendental $x=[0;a_1,a_2,\ldots]$ with
$a_n\in\{1,2\}$ and irrationality exponent $2$ for which
$7^{v_7(p_n^2-2q_n^2)}\ge q_n^\eta$ at infinitely many convergents,
for a fixed $\eta>0$. One can arrange that every finite word on
$\{1,2\}$ occurs.
\end{theorem}
\begin{proof}
Write $D_a=\left(\begin{smallmatrix}a&1\\1&0\end{smallmatrix}\right)$,
so $D_{a_1}\cdots D_{a_n}$ has first column $(q_n,p_n)^t$.
Use as generators the four products $D_aD_b$, $a,b\in\{1,2\}$.
Quotients obtained by varying the first or the last digit give the upper
and lower elementary shears of parameter $1$, up to sign. Thus the
symmetric generating set generates $\mathrm{SL}_2(\mathbb Z)$.
The expansion theorem of Bourgain and Varj\'u \cite[Thm.~1]{bv}
gives a common expansion constant in its reductions modulo $7^h$.
In a finite group the number of edges leaving a set $A$ under right
multiplication by $s$ equals the number leaving under $s^{-1}$.
Consequently positive word balls grow by a fixed factor up to half
the group. Forward and backward balls of radius $O(\log |G|)$, where
$G=\mathrm{SL}_2(\mathbb Z/7^h\mathbb Z)$, therefore intersect.
Every even prefix can reach any target modulo $7^h$ by appending
$O(h)$ digits in pairs.

Enumerate all words on the two digits. At each stage extend the existing
prefix by the next word and, if necessary, a digit $1$ to restore even
length $L$. Take $h=L+1$. Find an integer $r$ with
$v_7(r^2-2)=h$, by lifting $3$ modulo $7^h$ and choosing one of its
six lifts which fails modulo $7^{h+1}$. Append a shortest positive word
to reach $\left(\begin{smallmatrix}1&0\\r&1\end{smallmatrix}\right)$
modulo $7^{h+1}$. Exhaustive search terminates by the preceding paragraph.
The resulting convergent has $v_7(p_n^2-2q_n^2)=h$ and $n\le C h$
with one constant $C$. Since $q_n\le3^n$, this gives the required power
saving with $\eta=(\log 7)/(C\log 3)$. Every convergent has
$|x-p_n/q_n|<q_n^{-2}$, so Theorem~\ref{th:pelldrop} applies.
Bounded partial quotients give a uniform lower bound
$|x-a/b|\gg b^{-2}$ for all rationals and hence the asserted exponent.
\end{proof}

\subsection{A certified test on Ap\'ery's constant}\label{sec:aperytest}
Write $p_n/q_n$ for the convergents of $\zeta(3)$ itself, with $a_0=1$
and $p_0/q_0=1/1$. The first $1{,}000$ convergents give the following
distribution for $v_n=v_7(p_n^2-2q_n^2)$.
\[
\begin{array}{c|rrrrrr}
 v_n&0&1&2&3&4&\ge5\\ \hline
 \#\{0\le n<1000\}&738&230&24&7&1&0
\end{array}
\]
For example $p_1/q_1=5/4$ gives $p_1^2-2q_1^2=-7$.
The unique depth-four index is $n=692$. Its denominator has $364$
decimal digits, and $(p_{692},q_{692})\equiv(14647,13700)\pmod{16807}$, giving
$p_{692}^2-2q_{692}^2\equiv6\cdot7^4\pmod{7^5}$.
The local depth required by a fixed positive power saving must grow with
$\log q_n$ along infinitely many indices. This finite distribution leaves
the transcendence of $\zeta(3)$ open.

\paragraph{Certification.}
Use the inverse-binomial identity
$\zeta(3)=\sum_{j\ge1}(-1)^{j+1}t_j$, where
$t_j=5/(2j^3\binom{2j}{j})$
\cite[Eq.~(1.2)]{aperybb}. Set $N=2500$, $D=10^{1400}$ and
$W=\sum_{j=1}^N(-1)^{j+1}\lfloor Dt_j\rfloor$.
The $t_j$ decrease, and $Dt_{N+1}<1$ by integer comparison.
Rounding the $N/2$ positive and $N/2$ negative terms separately and
using the alternating-series remainder gives a rational enclosure
\[
 [L,U]=[(W-1251)/D,(W+1251)/D].
\]
To reproduce the table, initialize $\binom00=1$ and update
$\binom{2j}{j}=2(2j-1)\binom{2j-2}{j-1}/j$. On each continued-fraction
step, if $a=\lfloor L\rfloor=\lfloor U\rfloor<L$, record $a$ and
replace $[L,U]$ by $[1/(U-a),1/(L-a)]$. This procedure certifies
$1{,}388$ consecutive partial quotients, beginning $1,4,1,18,1,1,1,4$.
The usual numerator and denominator recurrences and repeated exact
division by $7$ then give the displayed counts.

\section{Ten quadratic-field lacunary quotients}\label{sec:tenpell}
Put $u=3+\sqrt2$ and $v=3-\sqrt2$. Thus $uv=7$ and $1<v<u$.
For an integer $k\ge2$ and positive integral weights $c_j$ of at most
exponential growth in $j$, set
\[
 G_{k,c}(z)=\sum_{j=0}^{\infty}c_jz^{k^j},\qquad
 X_{k,c}=\sqrt2\,
 \frac{G_{k,c}(1/v)+G_{k,c}(1/u)}{G_{k,c}(1/v)-G_{k,c}(1/u)} .
\]
Both sums converge. Positivity makes the denominator strictly positive.
The following ten choices include constant, geometric, polynomial and
second-order recurrence weights. The entries $F,L,P,J$ are the Fibonacci,
Lucas, Pell and Jacobsthal sequences, respectively. Their
initial data and recurrences follow the table.
\begin{center}\small
\begin{tabular}{@{}rccl r@{}}\toprule
Row & $k$ & Weight $c_j$ & $X_{k,c}$ (approx.) & $v_7(p^2-2q^2)$ at $t=2$\\\midrule
1 & 11 & $1$ & 2.97576725717231 & 121\\
2 & 13 & $13^j$ & 2.88222762964195 & 169\\
3 & 17 & $17^{2j}$ & 2.65115023905361 & 290\\
4 & 19 & $j+1$ & 2.99877018010686 & 361\\
5 & 23 & $\binom{j+2}{2}$ & 2.99970811554051 & 529\\
6 & 29 & $F_{j+1}$ & 2.99999388072752 & 841\\
7 & 31 & $L_j$ & 2.99999878330495 & 961\\
8 & 37 & $P_{j+1}$ & 2.99999969396311 & 1369\\
9 & 41 & $2^j$ & 2.99999995160561 & 1683\\
10 & 43 & $J_{j+1}$ & 2.99999999037777 & 1849\\
\bottomrule\end{tabular}\end{center}
$F$ and $L$ obey $s_{j+2}=s_{j+1}+s_j$ with initial pairs $(0,1)$
and $(2,1)$. For $P$ and $J$ the initial pair is $(0,1)$ and the
respective recurrences are $P_{j+2}=2P_{j+1}+P_j$ and
$J_{j+2}=J_{j+1}+2J_j$.

\begin{theorem}\label{th:pellten}
Every $X_{k,c}$ above with $v^k>7$ is transcendental by
Theorem~\ref{th:pelldrop}. If $v^k>49$, ordinary real approximation
already proves the assertion. These sufficient bounds hold for all
integers $k\ge5$ and $k\ge9$, respectively. Each of the ten entries
satisfies both bounds.
\end{theorem}
\begin{proof}
Truncate both sums at $j=t$ and denote the resulting quotient by $x_t$.
For $N=k^t$, write
\[
 A=7^N \sum_{j=0}^t c_jv^{-k^j}
   =u^N\sum_{j=0}^t c_jv^{N-k^j}
   =a+b\sqrt2\in\mathbb Z[\sqrt2].
\]
Its conjugate $A'=a-b\sqrt2$ is the corresponding sum at $1/u$.
Hence $x_t=a/b$. Write $x_t=p/q$ in lowest terms with $q>0$;
the convergent series give $b\ll 7^N$ and $q\ll 7^N$.
The ratio of the truncated sum at $1/v$ to that at $1/u$ is a
weighted average of $(u/v)^{k^j}$, with positive weights $c_ju^{-k^j}$.
It strictly increases with $t$, so $x_t$ strictly decreases to $X_{k,c}$.
In particular the reduced denominators tend to infinity. The differences of partial sums in the quotient defining $x_t$ stay bounded away from
zero. The sum of the two discarded tails is
$O(C^{t+2}v^{-kN})$ for an appropriate $C>1$, since consecutive
terms of the majorizing series have ratio at most $1/2$ for large $t$.
Subtracting the two quotients gives
\[
 0<x_t-X_{k,c}\le \exp\{-kN\log v+o(N)\}.
\]

At the two embeddings in $\mathbb Q_7$, the valuations of $u,v$ are
$(1,0)$ and $(0,1)$, respectively, because $uv=7$ and $u-v=2\sqrt2$
is a unit. At the first embedding $A$ has valuation at least $N$.
At the second, in the last expression for $A$ the final summand has
valuation $d=v_7(c_t)=O(t+1)$, whereas the other summands have
valuation at least $N-k^{t-1}$. Thus $v_7(A)=d$ at the second
embedding for large $t$. The linear transformation from $(a,b)$ to
the two embedded values has determinant a $7$-adic unit. Consequently
$v_7(\gcd(a,b))=d$, and taking the norm gives
\[
 v_7(p^2-2q^2)\ge N-d .
\]
The weighted expression of Theorem~\ref{th:pelldrop} is consequently
at most $\exp\{N(\log7-k\log v)+o(N)\}$. If $v^k>7$, choose
$0<\delta<k\log v/\log7-1$ and use $q\ll7^N$ to obtain a bound
$q^{-\delta}$ for large $t$. If $v^k>49$, the real error itself is
$<q^{-2-\varepsilon}$ for a fixed $\varepsilon>0$.
Reduced rational spacing excludes a rational limit with this accuracy
and Roth's theorem excludes an algebraic irrational limit.
For the ten rows there are infinitely many $t$ with $7\nmid c_t$, giving the sharper lower bound $N$ there. Indeed their companion matrices have determinant prime to seven, and each coefficient state returns periodically modulo seven to its initial state with $c_0\ne0\pmod7$. The last column of the table uses the reduced $p/q$ at $t=2$. Finally, $\log7/\log v$ lies between $4$ and $5$, and
$2\log7/\log v$ lies between $8$ and $9$.
\end{proof}


\subsection{Algebraic independence of the ten quotients}\label{sec:mahler}
The approximation bound has a numerical interpretation. At the smallest
tabulated base, $k=11$, the real exponent supplied by truncation is
$11\log v/\log 7=2.606433\ldots$ in the sense that every smaller exponent
is achieved eventually. For $5\le k\le8$ the same real estimate falls
below two, while Theorem~\ref{th:pelldrop} still applies.
A simultaneous assertion calls for information about relations between different rows.

\begin{theorem}\label{th:joint}
For the ten rows of Section~\ref{sec:tenpell}, the twenty numbers
$G_{k,c}(1/u)$ and $G_{k,c}(1/v)$ are algebraically independent over
$\overline{\Q}$. The ten numbers $X_{k,c}$ are algebraically independent as well.
\end{theorem}
\begin{proof}
Write $K=\Q(\sqrt2)$ and let $w=u$ or $v$. First,
$G_{k,c}(1/w)\notin K$ in every row. Indeed, suppose this value is
$b\in K$ and let $s_t$ be its partial sum through $t$, $N=k^t$.
Choose a positive integer $D$ for which $Db$ is an algebraic integer.
Then $Dw^N(b-s_t)$ is a nonzero algebraic integer. The conjugate of
$s_t$ is the partial sum at the other reciprocal and remains bounded.
Taking the absolute norm gives
\[
  1\le D^2 7^N |b-s_t|\,|b'-s_t'|\ll_b7^N |b-s_t|.
\]
The tail bound $|b-s_t|\le w^{-kN+o(N)}$ contradicts this inequality
since $w^k\ge v^k>7$.

Each table row satisfies a rational constant-coefficient recurrence with
nonzero trailing coefficient. Choose its order $d$ and set
$H_s(z)=\sum_{j\ge0}c_{j+s}z^{k^j}$ for $0\le s<d$.
With $H=(H_0,\ldots,H_{d-1})^t$, the identity
$H_s(z)=c_s z+H_{s+1}(z^k)$ and the recurrence give
$H(z)=z\boldsymbol c+C H(z^k)$ with a constant invertible companion
matrix $C$. Adjoining $1$ makes the system homogeneous with constant
nonzero determinant. The iterates of $G=H_0$ therefore lie in a finite-dimensional
$K(z)$-space, giving a scalar $k$-Mahler equation after clearing denominators.
All points of the open punctured unit disc are regular for this system.

The numbers $1/u$ and $1/v$ are multiplicatively independent. Their valuations
at the two primes of $K$ above $7$ give the vectors $(-1,0)$ and $(0,-1)$,
in a suitable order. The independent-points result of Adamczewski and
Faverjon \cite[Thm.~1.1(i)]{af} now gives algebraic independence of the two
values for each $k$, since neither belongs to $K$.

For each row take two copies of its system in independent variables, with
transformation $(z_1,z_2)\mapsto(z_1^k,z_2^k)$, at $(1/u,1/v)$.
The transformation matrix $kI_2$ has a positive eigenvector for its
spectral radius, both coordinates tend to zero along its orbit, and
multiplicative independence excludes a nontrivial monomial identically
one along any infinite arithmetic progression of iterates.
These are the admissibility conditions of \cite[\S5]{af}.
Regularity follows from the constant nonzero determinants.
The ten spectral radii are distinct primes. Purity for independent
transformations \cite[Thm.~3.8 and Cor.~3.9]{af} adds the ten
transcendence degrees, yielding $20$.

For each pair write $A=G(1/u)$ and $B=G(1/v)$. The substitution
$(A,B)\mapsto (A,X)$ is birational over $K$, since
$B=A(X+\sqrt2)/(X-\sqrt2)$ and $X>\sqrt2$. Applied independently in
the ten pairs, this substitution preserves transcendence degree $20$.
Its ten $X$-coordinates are consequently algebraically independent.
\end{proof}

For the constant weights $c_j=1$, existing Mahler theory also treats the
small bases $k=2,3,4$. To see this, write $G(z)=z+G(z^k)$.
The function has a singularity at every $k^r$-th root of unity. Indeed,
along a radius ending there the tail with $j\ge r$ is positive and
unbounded, while the finite initial sum stays bounded. An algebraic
function has only finitely many singular points in the finite plane,
so $G$ is transcendental over $\overline{\Q}(z)$.
At either reciprocal the regular system $(1,G)$ and the specialization
theorem \cite[Cor.~3.5]{af} give a transcendental value.
The independent-points result then gives two algebraically independent
values and hence a transcendental quotient for every $k\ge2$.


\section{A counting criterion for adjacent convergents}\label{sec:rankfour}
Write \(x=[0;a_1,a_2,\ldots]\) for an infinite regular continued fraction, and \(p_n/q_n\) for its reduced convergents. Set
\[
 M_n=\begin{pmatrix}q_n&q_{n-1}\\p_n&p_{n-1}\end{pmatrix},\qquad
 V_n=(q_n,p_n,q_{n-1},p_{n-1}).
\]
We use \(v_\ell(\ell)=1\) and \(v_\ell(0)=+\infty\). Each \(V_n\) is a primitive integral vector. Four such vectors are independent if their \(4\)-by-\(4\) determinant is nonzero.
\begin{theorem}\label{th:count}
Fix a prime \(\ell\). Suppose that for arbitrarily large integers \(m\) one can select \(m^4\) different indices with
\[
 q_n\ge e^{m^2},\qquad
 \ell^{v_\ell(p_n-q_{n-1})}\ge q_n^{1/m},
\]
such that any four selected vectors \(V_n\) are linearly independent over \(\Q\). Then \(x\) is transcendental.
\end{theorem}
\begin{proof}
Assume \(x\) is algebraic. Apply Evertse--Ferretti \cite[Thm.~3.1]{ef} over \(\Q\), in dimension four, at \(\infty\) and \(\ell\), with \(\varepsilon=1/m\). In the coordinates \((X_1,X_2,X_3,X_4)\) use
\[
\begin{array}{c|l|c}
v &\text{linear forms in the order used}&\text{normalized exponents}\\ \hline
\infty &X_1,\ X_2-xX_1,\ X_3,\ X_4-xX_3 &(0,-2,0,-2)\\
\ell &X_1,\ X_2,\ X_4,\ X_2-X_3 &(0,0,0,-1/m).
\end{array}
\]
Both determinants have absolute value one. At infinity \(\|V_n\|_\infty=q_n\), and both raw errors are \(<1/q_n\) by the convergent inequalities. At \(\ell\), primitivity gives \(\|V_n\|_\ell=1\). The height is \(q_n\), and the exponents sum to \(-4-1/m\). For large \(m\), \(e^{m^2}\) exceeds the theorem's threshold \(\max(C_x,4^{4m})\). Its bound for the number of proper rational subspaces is \(O_x(m^3\log^2(2m))\). Each such space contains at most three selected vectors. For sufficiently large \(m\), the resulting \(O_x(m^3\log^2(2m))\) bound for the number of selected vectors is smaller than \(m^4\).
\end{proof}


\subsection{A computable example with bounded partial quotients}
For a prime \(p\), write \(d_p\) for the projective distance on \(\PP^1(\Q_p)\). For vectors with maximum \(p\)-adic coordinate norm one, it is the \(p\)-adic absolute value of their determinant. We write \(a/b\) for \([a:b]\) in this notation.
\begin{theorem}\label{th:example}
There is a computable \(x=[0;a_1,a_2,\ldots]\) with the following properties.
\begin{enumerate}
\item Every \(a_n\) is \(7\) or \(14\), and every word in these two digits occurs. The irrationality exponent of \(x\) equals \(2\).
\item Theorem~\ref{th:count} applies at \(11\). For all sufficiently large \(m\), every selected index satisfies
\[
 \frac1m\le\frac{v_{11}(p_n-q_{n-1})\log 11}{\log q_n}\le\frac3m.
\]
\item Fix a finite set \(S\) of primes, algebraic targets \(\alpha_p\in\PP^1(\Q_p)\) for \(p\in S\), and \(\varepsilon>0\). Along every sequence of distinct reduced \(a/b\to x\), \(b>0\), one eventually has
\[
 b^2\,|x-a/b|\prod_{p\in S}d_p(a/b,\alpha_p)>b^{-\varepsilon}.
\]
In particular, for \(|x-a/b|\le C/b^2\) the logarithm of the reciprocal of the product of local distances is \(o(\log b)\).
\end{enumerate}
\end{theorem}
Here irrationality exponent means the supremum of the real \(\nu\) for which \(|x-a/b|<b^{-\nu}\) infinitely often. A bounded regular continued fraction satisfies \(|x-a/b|\gg_x b^{-2}\). Its convergents also give infinitely many upper bounds \(b^{-2}\).

\paragraph{Short continuations with a prescribed residue.}
Put \(D_a=\left(\begin{smallmatrix}a&1\\1&0\end{smallmatrix}\right)\), so \(M_n=D_{a_1}\cdots D_{a_n}\). Let \(S\) consist of the four products \(D_aD_b\) with \(a,b\in\{7,14\}\). If \((N,7)=1\), any even prefix can be extended to any prescribed matrix in \(\SL_2(\Z/N\Z)\) with \(O(\log(2N))\) digits.
Indeed, quotients of elements of \(S\) give upper and lower elementary shears with parameter \(7\), by varying the first or last digit. Their powers give all elementary shears modulo \(N\), which generate \(\SL_2(\Z/N\Z)\) by elimination over prime powers and the Chinese remainder theorem. Over \(\Z\) these shears generate a Zariski-dense subgroup. Bourgain--Varj\'u's theorem \cite[Thm.~1]{bv} gives uniform expansion for \(S\cup S^{-1}\) in all these quotients. For every subset \(A\) of the group and \(s\in S\),
\(\#(As\setminus A)=\#(As^{-1}\setminus A)\).
Thus forward balls of words of length at most the given radius in the positive generators grow by a fixed factor until they exceed half the group. Backward balls from the target do likewise, and the two balls meet. The group has size at most \(N^4\), giving the asserted length. A search for a shortest word terminates and implements this step.

\paragraph{Two simultaneous prescriptions.}
Start with \((a_1,a_2)=(7,7)\). At stage \(m\ge2\) take the least prime \(r_m>m^4\). Bertrand's postulate gives \(r_m<2m^4\). Number the \(m^4\) extensions in this stage by \(j=1,\ldots,m^4\). With \(Q\) the denominator before each extension, take the least \(k\) with \(11^{mk}\ge Q^2\). At its endpoint prescribe
\[
 M_n\equiv
 \begin{pmatrix}1&0\\11^k&1\end{pmatrix}T
 \pmod {11^{k+1}},\qquad
 M_n\equiv H(j):=
 \begin{pmatrix}1+j^3&j\\j^2&1\end{pmatrix}\pmod {r_m},
\]
where
\[
 t=11m,\qquad
 T=\begin{pmatrix}1+t&t^2\\ t^2&(1+t^4)/(1+t)\end{pmatrix}\in\SL_2(\Z_{11}).
\]
Both determinants equal one. The prescriptions can therefore be imposed after any even prefix. The off-diagonal difference has valuation exactly \(k\), since \(1+t\) is a unit. Every four endpoint vectors of a stage are independent. Indeed, their reductions modulo \(r_m\) are evaluations of \(1,j,j^2,j^3\) up to an invertible linear change of coordinates. Their determinant is a signed Vandermonde product and \(r_m>m^4\).

The entries of the unperturbed matrix \(T(11m)\) have rational height polynomial in \(m\). We use this bound for fixed algebraic affine forms on the symmetric matrices. Indeed, the rational functions
\[
 1,\quad 1+t,\quad t^2,\quad (1+t^4)/(1+t)
\]
are linearly independent over the algebraic constants, by coefficient comparison after multiplication by \(1+t\). For any fixed algebraic affine form \(F\) on four entries that is not proportional to the off-diagonal difference, \((1+t)F(T(t))\) is thus a nonzero polynomial of degree at most four. Clear a fixed denominator so its coefficients are algebraic integers in a field of degree \(d\). At \(t=11m\) its norm is a nonzero integer for all sufficiently large \(m\), of ordinary size \(O_F((1+m)^{4d})\). Thus at the chosen place over \(11\), \(v_{11}(F(T(11m)))=O_F(\log m)\) in absolute value, since \(1+11m\) is a unit. The congruence \(M_n\equiv T(11m)\bmod 11^k\) changes the value of \(F\) by a quantity of valuation at least \(k-O_F(1)\). Below we obtain \(k\gg m^4\), so these valuations agree eventually. For a nonzero multiple of the off-diagonal difference the valuation is exactly \(k+O_F(1)\). Neither type of form gives a fixed positive exponent on the selected endpoints.
\subsection{Avoiding simultaneous approximation to fixed local algebraic points}

Use the projective coordinates $a/b=[a:b]$ for a reduced rational
with $b>0$. Fix a finite set $S$ of primes and algebraic points
$\beta_\ell\in\PP^1(\Q_\ell)$ for $\ell\in S$. Choose algebraic coordinates in a local chart with one coordinate equal to one and the other integral. Put
\[
 d_\ell([a:b],[u:v])=|av-bu|_\ell,\qquad
 V_\ell=d_\ell(a/b,\beta_\ell)^{-1},\qquad V=\prod_{\ell\in S}V_\ell .
\]
The rational vector $(a,b)$ has local sup norm one at every prime.
Each finite $V_\ell$ is a nonnegative integral power of $\ell$.
An exact equality with a rational target singles out one rational
point and can be omitted when $b\to\infty$.

\begin{lemma}\label{lem:weighted}
The choices in the matrix construction can be made computably so that
every finite word over $\{7,14\}$ occurs and, for every such fixed tuple
and every $e>0$, there is no sequence of distinct reduced $a/b\to x$
satisfying
\begin{equation}\label{eq:weighted-bad}
 b^2|x-a/b|\,V^{-1}\le b^{-e}.
\end{equation}
\end{lemma}
\begin{proof}
For a fixed tuple, any bounded sequence of distinct rational slopes
satisfies $V\le b^{2+\eta}$ eventually, for every fixed $\eta>0$.
This follows from the projective subspace theorem, for instance
\cite[Cor.~3.2]{ef}. For the corresponding Roth--Ridout formulation see \cite[Thm.~1]{cor}. At each $\ell$ use the local determinant
form and a coordinate form completing it to a system of determinant
of local norm one. Its second factor is at most one. The normalized
product is thus at most $V^{-1}$. The global projective height
$\max(|a|,b)$ is comparable to $b$. Infinitely many exceptions would
satisfy the theorem's inequality with excess $\eta/2$ and lie on
finitely many rational lines. Each line supplies just one rational
slope. This argument applies also when a target is infinity. For
$S=\varnothing$, bounded partial quotients already give the assertion.

For $S\ne\varnothing$ and $0<e\le1$ define the positive real quantity
\[
 H_e= b^{1+e/4}V^{-1/2}.
\]
The real neighborhoods in \eqref{eq:weighted-bad} have radius
$R=Vb^{-2-e}=b^{-e/2}H_e^{-2}$. We choose our intervals to control
them at the scale $H_e$, even when that scale is much smaller than $b$.

Enumerate the tests comprising a finite nonempty tuple of distinct
primes with algebraic projective targets and an exponent $e=1/j$,
$j\ge1$. Integer squarefree polynomials with Hensel-isolating discs
in the two integral charts enumerate all local targets. The record
with $s=|S|$ and $e=1/j$ is assigned integer cost
$w=100j(2+128sj)^s$. At stage $m$, activate the longest initial
segment of records whose total cost is at most $\sqrt m$ and whose
indices $j$ are all at most $\log_2(m+2)$. Each fixed record
eventually stays active.

Each step starts with the even parent prefix $G=M_n$ of denominator
$Q=q_n$. Let $h=\lceil\log_2(m+1)\rceil$. Prepare $2^h$ candidate
continuations using $h$ successive choices between $(7,7)$ and
$(14,14)$. At the first step of this stage insert word number $m-1$ in length-lexicographic order, padding with one $7$ if needed
for parity. With the empty word indexed zero, these insertions contain
every nonempty word as a factor. Complete each candidate by a
shortest word in digit pairs achieving the two prescribed congruences
for the current index and $k$. Break ties lexicographically.
Writing $Q_*$ for the maximum candidate endpoint denominator, the
routing bound gives
\begin{equation}\label{eq:guard-growth}
 \log(Q_*/Q)\ll k+\log m\ll(\log Q)/m+\log m .
\end{equation}
Before stage $m$ at least $\sum_{i=2}^{m-1}i^4$ digit pairs have been
appended. Thus $\log Q\gg m^5$, and $Q_*/Q\le Q^{C/m}$ with an
absolute constant $C$ at all sufficiently large stages.

Let $I$ be the closed parent cylinder, of length at most $Q^{-2}$.
For each active record, list the fractions satisfying
\[
 Q^{1/2}\le b\le Q^m,\qquad V_\ell\le b^3\ (\ell\in S),\qquad
 Q\le H_e\le Q_*,\qquad
 [a/b-R,a/b+R]\cap I\ne\varnothing .
\]
All centers lie in an interval of length at most $3Q^{-2}$, since
$R\le Q^{-2}$. This bounds the numerator for each allowed denominator.
Equality with a rational target is decidable using its defining
polynomial and isolating disc. At other points successive local
digits determine each finite $V_\ell$ exactly. Endpoints of $I$ are
rational, and all remaining comparisons can be performed by taking
integer powers.

To count the listed fractions for one record, split the denominator
range into bins $Y\le b<Y^{1+e/16}$, starting at $Y=Q^{1/2}$.
There are $O(e^{-1}\log(2m))$ bins. Within a bin,
$0\le\log_Y V_\ell<4$. Partition each such interval into cells of
width $e/(32s)$, with at most $2+128s/e$ choices per prime.
Two distinct fractions in the same denominator bin and the same
local cells have a nonzero integral cross determinant $D$. With
subscripts $1,2$ indicating the fractions, we have
\[
 |D|\le 3Y^{2+e/8}Q^{-2},\qquad
 \prod_{\ell\in S}\min(V_{1,\ell},V_{2,\ell})
 \ge V_1Y^{-e/32}\ge Y^{2+e/2-e/32}Q_*^{-2}.
\]
At each prime, normalized primitive vectors at projective distance
at most $\ell^{-a}$ from one target are proportional modulo
$\ell^a$ by unit scalars. The product on the left therefore divides
$D$. For all active records at sufficiently large stages,
$(Q_*/Q)^2\le Y^{e/16}$, by \eqref{eq:guard-growth},
$Y\ge Q^{1/2}$ and $e\ge1/\log_2(m+2)$. The lower bound then
exceeds the upper bound, uniformly, since $Y^{e/4}\to\infty$.
Consequently each joint cell contains at most one center. The
cost budget bounds the total number of marked neighborhoods by
$O(\sqrt m\log(2m))$.

Different guard cylinders have gap at least
$(1/8-1/14)/(4\cdot15^{4h}Q^2)$. At the first differing pair their
next tail coordinates lie respectively in $[1/8,1/7]$ and
$[1/15,1/14]$. The common-prefix map has absolute derivative at least
$1/(4Q_{\rm pre}^2)$, with $Q_{\rm pre}\le 15^{2h}Q$.
Each marked neighborhood has diameter at most
$2Q^{-2-e/4}$, smaller than this gap for all active records once
$m$ is large. Hence some guard survives among the $2^h>m$ candidates.
Commit to the lexicographically first survivor and its completion.
If no guard survives, commit to the first candidate.
The estimates show that this fallback occurs at only finitely many
stages. This rule avoids any need to compute asymptotic constants.

Finally fix an exponent $e=1/j$ and its local tuple. For a proposed
infinite sequence in \eqref{eq:weighted-bad}, use the finite-place
bound with $\eta=e/4$. Eventually $V_\ell\le V\le b^{2+e/4}$ and
$b^{e/8}\le H_e\le b^{1+e/4}$. Thus $H_e\to\infty$. Successive
committed denominators $Q,Q'$ cover every large value of $H_e$ by
intervals $[Q,Q']$. In its covering step the record is eventually
active, $Q'\le Q_*\le Q^2$, and
$Q^{1/2}\le b\le Q^{16/e}\le Q^m$. The center would therefore have
been listed and its closed neighborhood avoided by the committed
guard, contrary to \eqref{eq:weighted-bad}. Every fixed positive
exponent dominates some $1/j$, proving the lemma.
\end{proof}
\begin{proof}[Completion of the proof of Theorem~\ref{th:example}]
Commit to the choices of Lemma~\ref{lem:weighted}. For every stage-\(m\) endpoint with denominator \(Q'\), one has \(\log Q\le\log Q'\le(1+C/m)\log Q\) for large \(m\). By definition,
\[
 \frac{2\log Q}{m}\le k\log 11<\frac{2\log Q}{m}+\log 11 .
\]
Since \(\log Q\gg m^5\), we obtain \(1/m\le k\log11/\log Q'\le 3/m\) and \(Q'\ge e^{m^2}\) eventually. The modular Vandermonde determinant proves the remaining hypothesis of Theorem~\ref{th:count}. Infinite bounded digits give irrationality exponent two, and the inserted words give exactly \(2^d\) distinct blocks of each length \(d\ge1\). Lemma~\ref{lem:weighted} supplies the weighted inequality. If \(b^2|x-a/b|\le C\), the same inequality gives \(1\le V<C b^\varepsilon\) eventually for every \(\varepsilon>0\), proving the last assertion.
\end{proof}

\subsection{Comparison with the split quadratic norm}
For the two projective roots of \(X^2-2Y^2\) over \(\Q_7\), a reduced pair \((a,b)\) with \(7\mid a^2-2b^2\) has \(7\nmid b\). Exactly one of \(a-b\sqrt2\), \(a+b\sqrt2\) is divisible by \(7\). It follows from Theorem~\ref{th:example}(3), applied separately to the two roots, that for every \(\varepsilon>0\) and every sequence of distinct reduced \(a/b\to x\), eventually
\[
 b^2|x-a/b|\,7^{-v_7(a^2-2b^2)}>b^{-\varepsilon}.
\]
When \(7\nmid a^2-2b^2\) the same conclusion follows from bounded partial quotients. Thus the comparison covers even the weighted form of the one-fraction Pell test. On the ordinary convergents there is the simpler identity \(M_n\equiv\left(\begin{smallmatrix}0&1\\1&0\end{smallmatrix}\right)^n\bmod7\), which gives \(7\nmid p_n^2-2q_n^2\) for every \(n\).

Earlier work studies shrinking exponents and the density of rational approximants \cite{bug2007}, as well as continued-fraction symmetries and transcendence \cite{ab}. The criterion here specializes a published quantitative theorem. Its comparison with one-fraction tests holds for fixed finite sets of primes, fixed algebraic targets, and fixed positive power savings. Quantitative counting results for one-fraction tests, number-field approximants and other transcendence methods remain separate possibilities. Application to a prescribed constant requires a proof of the congruence count and the four-point independence for that constant's actual convergents.

\section{Wolf's prime continued fraction}\label{sec:wolf}
Wolf \cite{wolf} proposed the transcendence of \(u=[0;2,3,5,7,\ldots]\), whose partial quotients are the primes in increasing order. Unbounded, strictly increasing partial quotients exclude quadratic irrationality by Lagrange's theorem. Its irrationality exponent equals two. For the $n$th prime $a_n$, Bertrand's postulate gives $a_n\le 2^n$, whereas $q_n\ge (n+1)!$. The usual continued-fraction formula $\mu(u)=2+\limsup_n\log a_{n+1}/\log q_n$ gives the claim. To go further, write \(p_n/q_n\) for the convergents of an infinite regular continued fraction \(x=[0;a_1,a_2,\ldots]\) and put
\[
V_n=(q_n,p_n,q_{n-1},p_{n-1}).
\]

\begin{lemma}\label{lem:four}
If the positive integer partial quotients are eventually strictly increasing, then all sufficiently late indices \(i<j<k<l\) with \(k-j\ge3\) give four linearly independent vectors \(V_i,V_j,V_k,V_l\) over \(\Q\).
\end{lemma}
\begin{proof}
Let \(r_n=q_{n-1}/q_n\), \(h_n=[0;a_{n+1},a_{n+2},\ldots]\), and \(t_n=(p_{n-1}-xq_{n-1})/q_n\). The identity \(x=(p_n+h_np_{n-1})/(q_n+h_nq_{n-1})\) gives
\[
p_n-xq_n=-h_n(p_{n-1}-xq_{n-1}),\qquad
(2q_n^2)^{-1}<|t_n|<q_n^{-2}.
\]
Consequently, dividing each row by its denominator and applying the same invertible column transformation gives rows \((1,r_n,t_n,t_nh_n)\). Eventually \(r_n,h_n\) decrease strictly and \(|t_n|\) decreases. Indeed \(r_n=1/(a_n+r_{n-1})\), \(h_n=1/(a_{n+1}+h_{n+1})\), with the remainders below \(1/2\), and \(q_{n+1}\ge2q_n\). For late indices this gives
\[
h_i-h_j\ge\frac1{2a_{j+1}^2},\qquad
r_k-r_l\ge\frac1{8a_k^2}.
\]
For the first bound use \(a_{i+1}\le a_{j+1}-1\). For the second compare \(1/(a_k+1/2)\) with \(1/(a_k+1)\).

Expand the determinant by its last two columns. The term using rows \(i,j\) in these columns has magnitude at least
\(\lvert t_i t_j\rvert/(16a_{j+1}^2a_k^2)\). Each of the five other terms has magnitude at most \(\lvert t_i t_k\rvert\), since all \(r,h\) lie between zero and one. The ratio of their combined magnitudes to the first term is at most
\[
160a_{j+1}^2a_k^2(q_j/q_k)^2\le160/a_{j+2}^2,
\]
using \(q_k/q_j\ge a_{j+1}a_{j+2}a_k\). This is eventually less than one.
\end{proof}

\begin{corollary}\label{cor:wolf}
For the prime continued fraction \(u\), fix a prime \(\ell\). It suffices for transcendence that, for arbitrarily large integers \(m\), at least \(3m^4\) distinct indices satisfy
\[
q_n\ge e^{m^2},\qquad \ell^{v_\ell(p_n-q_{n-1})}\ge q_n^{1/m}.
\]
\end{corollary}
\begin{proof}
One residue class modulo three supplies \(m^4\) indices. These are eventually all in the range of the lemma, so any four are independent. Theorem~\ref{th:count} now yields the conclusion.
\end{proof}
The lemma settles the independence requirement for the actual prime continued fraction. The needed lower bound for its congruence counts remains unproved.

\paragraph{A finite test at \(\ell=11\).}
Generate the first \(100{,}000\) primes by a sieve through \(1{,}299{,}709\). Start the continuant matrix at the identity and right-multiply by
\(\left(\begin{smallmatrix}a_n&1\\1&0\end{smallmatrix}\right)\) modulo \(11^6\).
Omit \(n=1\), where the off-diagonal difference is zero and the denominator is two. The counts for the actual prime continued fraction are
\[
\begin{array}{c|rrrrrr}
v_{11}(p_n-q_{n-1}) &0&1&2&3&4&\ge5\\ \hline
\#\{2\le n\le100{,}000\} &90{,}852&8{,}282&787&72&6&0 .
\end{array}
\]
The depth-four indices are \(9014,26483,46137,72187,82821,83037\). None meets the two displayed inequalities of the theorem for any \(m\ge2\). Indeed, writing \(v=v_{11}(p_n-q_{n-1})\), eligibility implies \(\log q_n\le v^2(\log 11)^2\). In the same recurrence, all indices with \(v>0\) have \(n\ge17v^2\) except \(n=6,v=1\). Use \(q_n\ge2^n\) and \(17\log2>(\log11)^2\), and check \(q_6>2^9\) directly. This finite exclusion gives no upper bound beyond the tested prefix.
\subsection{An obstruction for factorial partial quotients}
Let \(y=[0;1!,2!,3!,\ldots]\). At any fixed prime \(\ell\), the reduction
of its continuant matrices modulo \(\ell^2\) becomes periodic with period
dividing two from index \(2\ell-1\) onward. Indeed, \(\ell^2\mid n!\)
for \(n\ge2\ell\), and every subsequent multiplication is by the exchange
matrix \(\left(\begin{smallmatrix}0&1\\1&0\end{smallmatrix}\right)\)
modulo \(\ell^2\).
Write the matrix at index \(2\ell-1\) as
\(\left(\begin{smallmatrix}A&B\\C&D\end{smallmatrix}\right)\).
The off-diagonal differences \(p_n-q_{n-1}\) are therefore congruent to
\(C-B\) at odd indices and to \(D-A\) at even indices on this tail.

The exact computation for the primes \(\ell\le10{,}000\) gives the following
stabilized valuations.
\[
\begin{array}{c|c|l}
 (v_{\rm odd},v_{\rm even}) & \text{number of primes} & \text{primes}\\ \hline
 (1,0)&3&2,29,47\\
 (0,1)&2&5,19\\
 (0,0)&1224&\text{all other primes in the range}
\end{array}
\]
The following code reproduces the table using only integer arithmetic.
\begin{quote}\small
\begin{verbatim}
from math import isqrt
for ell in range(2, 10001):
    if any(ell % j == 0 for j in range(2, isqrt(ell) + 1)):
        continue
    M = ell**2
    A, B, C, D, f = 1, 0, 0, 1, 1
    for n in range(1, 2*ell):
        f = f*n % M
        A, B, C, D = (A*f+B) % M, A, (C*f+D) % M, C
    odd, even = (C-B) % M, (D-A) % M
    assert odd and even
    if odd % ell == 0 or even % ell == 0:
        print(ell, int(odd % ell == 0), int(even % ell == 0))
\end{verbatim}
\end{quote}
For any fixed prime in this range, the two size conditions of the
four-dimensional criterion force
\(m^2\le\log q_n\le m v_\ell(p_n-q_{n-1})\log\ell\).
On its eventual tail this forces \(m\le\log\ell\). The finite initial
segment also drops out for large \(m\) by the height condition, even
if a difference there vanishes. Thus each of the tested primes is
excluded as a witness for this criterion on \(y\). Larger primes are
outside the computation.



\section{Rotation codings of modular-form coefficients}\label{sec:hecke}
For an integer \(a\) let \([a]_m\) be its least nonnegative residue. Fix integers \(B,m,k\ge2\). For an integer sequence \(f=(f(n))_{n\ge1}\) define
\[
 \theta_{f,B,m}=\sum_{n\ge1}\frac{[f(n)]_m}{B^n},\qquad
 d_n=1+[\lfloor n\theta_{f,B,m}\rfloor]_k,\qquad
 X(f;B,m,k)=[0;d_1,d_2,\ldots].
\]
Ramanujan's function \(\tau\) is given by \(\sum_{n\ge1}\tau(n)z^n=z\prod_{n\ge1}(1-z^n)^{24}\).
Chatterjee and Mandal \cite[Thm.~1.11(1)]{cm} proved transcendence for \(f=\tau,\ B=10,\ m\in\{5,7,8,9\}\), and every \(k\ge2\). Their Remark~3.1 identifies a limitation of their criterion at \(m=3,4\). Theorem~1.1 of their paper separately treats nonperiodicity of coefficients at \(691\).
\begin{theorem}\label{th:tau}
For every \(B,m,k\ge2\), \(\theta_{\tau,B,m}\) is irrational and \(X(\tau;B,m,k)\) is transcendental. Its irrationality exponent is two.
\end{theorem}

\begin{lemma}\label{lem:zero}
Let the nonnegative integer coefficients \(a_n\) be bounded above by \(M\). Suppose they contain arbitrarily long strings of zeros and are positive infinitely often. Then \(\sum_{n\ge1}a_n B^{-n}\) is irrational for every integer \(B\ge2\).
\end{lemma}
\begin{proof}
If the string of length \(J\) starts at \(N+1\), multiply the sum by \(B^N\) and subtract the integer supplied by its first \(N\) terms. The remainder \(R_N\) satisfies
\[
 0<R_N\le \frac{M}{B-1}B^{-J}.
\]
For a sum of denominator \(v\), each positive remainder is at least \(1/v\). Taking \(J\) arbitrarily large gives a contradiction. The argument also covers coefficients larger than the base.
\end{proof}
\begin{proof}[Proof of the Ramanujan classification]
For a prime \(p\nmid m\) put \(u_e=\tau(p^e)\) for \(e\ge0\) and \(u_{-1}=0\). The Hecke recurrence and multiplicativity at coprime integers are the classical identities of Mordell \cite{mordell}. Write
\[
 A_p=\begin{pmatrix}\tau(p)&-p^{11}\\1&0\end{pmatrix},\qquad
 A_p^e\binom10=\binom{u_e}{u_{e-1}}\quad(e\ge0).
\]
The determinant is a unit modulo \(m\). Thus the reduction of \(A_p\) has finite order \(h_p\ge2\), and
\[
 \tau(p^{h_p-1})\equiv0,\qquad \tau(p^{jh_p})\equiv1\pmod m\quad(j\ge1).
\]
To produce a zero string, take \(J\) distinct primes \(p_i\nmid m\), put \(s_i=h_{p_i}-1\), and require
\[
 N+i\equiv p_i^{s_i}\pmod {p_i^{s_i+1}}\qquad(1\le i\le J).
\]
The Chinese remainder theorem gives arbitrarily large such \(N\). The exponent of \(p_i\) in \(N+i\) is exactly \(s_i\), hence \(\tau(N+i)\equiv0\pmod m\) by multiplicativity. The residues \(1\) at powers of a single prime supply the other hypothesis of Lemma~\ref{lem:zero}. We have proved irrationality of \(\theta=\theta_{\tau,B,m}\).

The rotation criterion for irrational $0<\theta<1$ goes back to Adamczewski, Bugeaud and Davison \cite[Thm.~4.1]{abd}. For our unrestricted slope, set \(\omega=\theta/k\). The digits are \(1+\lfloor k\{n\omega\}\rfloor\). They record the orbit of an irrational rotation through the \(k\) half-open arcs of length \(1/k\). An eventual period \(T\) would make the first arc invariant, up to its endpoints, under translation by \(T\omega\), since otherwise the symmetric difference contains an open interval visited by the dense tail orbit. The first arc has no nonzero translation with this property. For a block of length \(d\), the preimages of the \(k\) partition boundaries under the first \(d\) iterates supply at most \(kd\) endpoints. The word is constant on each resulting half-open arc. Thus its block complexity is at most \(kd\). Bugeaud's theorem \cite[Thm.~1.1]{bug2013} rules out algebraicity for a non-eventually-periodic continued fraction of this complexity. Bounded partial quotients give irrationality exponent two.
\end{proof}

\paragraph{Specific values.}
At \(B=10,\ k=2\), two instances outside the range of \cite[Thm.~1.11(1)]{cm} are
\[
 X(\tau;10,3,2)=0.61807280254448190389\ldots,\quad
 X(\tau;10,691,2)=0.36602550839044533867\ldots .
\]
For instance, the first 115 coefficients can be computed as a truncated integer product for \(\Delta\). For either modulus, enclose the slope between the partial sum and that sum plus the tail bound. For each of the first 111 positive multiples the two endpoints have the same floor. The convergents defining this 111-digit cylinder have endpoints with identical displayed decimals. The tail bound used here is $(m-1)/(9\cdot10^{115})$. For the positive-weight holomorphic modular forms considered here, arbitrarily long strings of zero residues also follow from Serre's density-one divisibility theorem \cite{serre}. 

\subsection{Five Eisenstein series and a quadratic boundary case}
The positive-index Fourier coefficients of \(E_4,E_6,E_8,E_{10},E_{14}\) have the respective forms \(c\,\sigma_s(n)\) for
\[
 (c,s)=(240,3),\ (-504,5),\ (480,7),\ (-264,9),\ (-24,13).
\]
Here \(\sigma_s(n)=\sum_{d\mid n}d^s\). These are families (2)--(6) in \cite[Thm.~1.11]{cm}. For all \(B,m,k\ge2\) one has the exact classification
\[
 X(c\,\sigma_s;B,m,k)=\frac{\sqrt5-1}{2}\quad(m\mid c),\qquad
 X(c\,\sigma_s;B,m,k)\text{ is transcendental}\quad(m\nmid c).
\]
Indeed, \(m\mid c\) makes every partial quotient one. Otherwise put \(m'=m/\gcd(m,c)>1\). The multiplicative function \(\sigma_s\) satisfies the same matrix argument modulo \(m'\), with trace \(1+p^s\) and determinant \(p^s\). This gives arbitrarily long runs with \(m'\mid\sigma_s(n)\), and infinitely many values congruent to one modulo \(m'\). After multiplication by \(c\) the residues modulo \(m\) therefore meet both hypotheses of Lemma~\ref{lem:zero}. The rotation proof applies verbatim.

In the Nathanson family of \cite[Thm.~1.11(7)]{cm}, the modulus \(3\) gives a different answer. The count \(\Phi(n)\) of nonempty subsets of \(\{1,\dots,n\}\) whose gcd is coprime to \(n\) satisfies \(\Phi(1)=1,\ \Phi(2)=2\). For \(n>1\), Möbius inversion gives \(\Phi(n)=\sum_{d\mid n}\mu(d)2^{n/d}\). Modulo three the function \(2^u\) equals \(-1+2\mathbf 1_{2\mid u}\), so convolution with \(\mu\) gives zero for \(n>2\), in agreement with \cite{tang}. For the base-ten coding at \(m=3\), the extra final reduction modulo ten in \cite{cm} does nothing. The slope is exactly \(1/10+2/100=3/25\). Its rotation coding is periodic (a period is \(25k\)), and its infinite regular continued fraction is a quadratic irrational.

\subsection{The inner slope at modulus two and a second rational slope}

The inner slope for Ramanujan's function at modulus two has a stronger classification.
For formal power series modulo two, Jacobi's identity
\[
 \prod_{r\ge1}(1-z^r)^3
   =\sum_{j\ge0}(-1)^j(2j+1)z^{j(j+1)/2}
\]
and the eighth-power Frobenius map show that
\(\Delta(z)\equiv\sum_{j\ge0}z^{(2j+1)^2}\pmod2\).
Consequently, with \(q=1/B\),
\[
\theta_{\tau,B,2}=\sum_{j\ge0}q^{(2j+1)^2}=\tfrac12\vartheta_2(q^4),
\qquad
\vartheta_2(t)=2\sum_{j\ge0} t^{(j+1/2)^2}.
\]
The theta-value theorem of Duverney, Ke.~Nishioka, Ku.~Nishioka and Shiokawa
\cite{dnns-theta} gives transcendence at every such algebraic nome. The
general Hecke argument establishes irrationality of the inner slope for
each modulus.

Let \(g(n)\) count the nonempty subsets of \(\{1,\ldots,n\}\) of gcd one,
as in \cite{cm}. For \(n\ge2\) the new subsets counted by
\(g(n)-g(n-1)\) contain \(n\). Among the subsets counted by \(\Phi(n)\)
the operation of inserting or removing \(n\) pairs these sets with
sets that omit \(n\). No singleton \(\{n\}\) qualifies, so no pair
contains the empty set. Hence \(2(g(n)-g(n-1))=\Phi(n)\).
The residues of \(\Phi(n)\) modulo three give
\(g(1)\equiv1\) and \(g(n)\equiv2\pmod3\) for every \(n\ge2\).
In any base \(B\ge2\) this yields
\[
 \theta_{g,B,3}=\frac{B+1}{B(B-1)},\qquad
 \theta_{\Phi,B,3}=\frac{B+2}{B^2}.
\]
For \(B=10\) these slopes are \(11/90\) and \(3/25\), respectively.
For a rational slope \(a/b\) with positive denominator, translation
of the index by \(bk\) preserves every coded digit. Both outer
continued fractions are therefore quadratic irrationals for every \(k\ge2\).



\section{Shifted products of odd-indexed Fibonacci numbers}
Let \(F_0=0,F_1=1,F_{n+2}=F_{n+1}+F_n\), and
\(S=\sum_{n\ge0}1/F_{2n+1}\). The transcendence of \(S\) follows from
the stronger algebraic independence theorem of Elsner, Shimomura and
Shiokawa \cite{ess-odd}. For an odd number \(d=2h+1\) of distinct
nonnegative shifts \(r_1,\ldots,r_d\), put
\[
 T(\boldsymbol r)=\sum_{n\ge0}\prod_{j=1}^d \frac{1}{F_{2(n+r_j)+1}} .
\]
\begin{theorem}
There are real algebraic \(C>0\) and \(D\) such that
\(T(\boldsymbol r)=CS+D\). In particular every such shifted product
sum is transcendental.
\end{theorem}
\begin{proof}
Put \(\phi=(1+\sqrt5)/2\), \(t_j=\phi^{4r_j}\), and
\(z=\phi^{2(2n+1)}\). Binet's formula gives
\(1/F_{2(n+r_j)+1}=\sqrt{5t_jz}/(1+t_jz)\).
The simple-pole decomposition of
\(z^h/\prod_j(1+t_jz)\) thus gives
\[
 \prod_j\frac{1}{F_{2(n+r_j)+1}}
  =\sum_j \frac{c_j}{F_{2(n+r_j)+1}},
 \qquad
 c_j=K(-1)^h\frac{t_j^{h-1/2}}{\prod_{i\ne j}(t_j-t_i)},
 \quad K=5^h\prod_i\sqrt{t_i}.
\]
Indeed multiplication by \(1+t_jz\) and substitution of \(z=-1/t_j\)
determine each coefficient. The sum \(C=\sum_j c_j\) equals \(K(-1)^h\)
times the divided difference of \(t^{h-1/2}\) on the \(d\) distinct
positive nodes. When \(h\ge1\), the mean-value theorem for divided
differences evaluates it as the \(2h\)-th derivative at an intermediate
point divided by \((2h)!\). Exactly \(h\) factors of that derivative's
coefficient are negative, so \(C>0\). For \(h=0\) it equals one.
Summing the termwise identity, which involves absolutely convergent
series, gives \(D=-\sum_j c_j\sum_{n<r_j,\,n\ge0}1/F_{2n+1}\).
\end{proof}
For example \(T(0,1,2)=S/3-1/2\), since the three coefficients are
\((-1/3,1,-1/3)\). The full reciprocal sum in the original catalogue is
\(\sum_{n\ge1}1/F_n\).


\appendix
\section{The two selections from the original catalogue}\label{sec:rankings}
The first review ranked entries by their mathematical interest across the
full catalogue. A second selection restricted attention to the entries
whose recorded conclusion was irrationality. The two columns reproduce
those judgments in their original order. Each identifier locates the
definition and proof in the companion catalogue \cite{catalogue}.

\begin{center}
\small\setlength{\tabcolsep}{3pt}\renewcommand{\arraystretch}{1.25}
\begin{tabularx}{\linewidth}{@{}r l >{\raggedright\arraybackslash}X l >{\raggedright\arraybackslash}X@{}}
\toprule
Rank & ID & All catalogue entries & ID & Irrationality selection\\
\midrule
1 & M19 & Reflected $q$-Stieltjes zero $Q_*$ & A22 & Teardrop elastica arc-length ratio\\
2 & A16 & Half-energy radius of an annular aperture & A13 & Optimal annular Steklov eigenvalue\\
3 & A22 & Teardrop elastica arc-length ratio & M24 & Gieseking cusp packing density\\
4 & M16 & Generalized pentagonal composition growth & M25 & Zeta derivative $\zeta'(-2)$\\
5 & A08 & First upper-half-plane zero of $\operatorname{erfc}$ & N01 & Plastic-base value $\gamma_0(\rho,1)$\\
6 & M03 & Negative second derivative at the modular-discriminant peak & N12 & Ap\'ery's constant $\zeta(3)$\\
7 & M05 & Simple-cubic lattice resistance to $(1,1,0)$ & N07 & Reciprocal Fibonacci sum\\
8 & A09 & Location of a modified-Bessel maximum & N23 & Prime-square reciprocal Fibonacci sum\\
9 & A13 & Optimal annular Steklov eigenvalue & N18 & Binary coprime square-shell sum\\
10 & G01 & Random-graph $3$-core threshold & N20 & Tschakaloff value $T_\rho(1)$\\
\bottomrule
\end{tabularx}
\end{center}
The two labels $\gamma_0(\rho,1)$ and $T_\rho(1)$ use the plastic
number $\rho$, the real root of $x^3-x-1$. Euler's constant $\gamma$
is a distinct target. For $\gamma$, a one-zero argument still requires a
specific transcendental $T$ and a verified recovery identity
$\gamma=R(T)$. For the convergent-pair approach it would instead require
the two quantitative hypotheses of Theorem~\ref{th:count} on $\gamma$'s
own convergents, under an assumption of algebraicity. The constructions
above provide neither input for this particular constant.


\clearpage\phantomsection\addcontentsline{toc}{section}{References}
\begin{thebibliography}{99}\small

\bibitem{catalogue} \emph{One-zero proofs of irrationality and transcendence. A catalogue of one hundred arithmetic conclusions}, companion manuscript (15 November 2026). Entries are cited by their catalogue identifiers.

\bibitem{abd} B. Adamczewski, Y. Bugeaud, and L. J. Davison, \emph{Continued fractions and transcendental numbers}, Ann. Inst. Fourier \textbf{56} (2006), 2093--2113. \href{https://arxiv.org/abs/math/0511682}{Author version}.

\bibitem{af} B. Adamczewski and C. Faverjon, \emph{Mahler's method in several variables and finite automata}, Ann. of Math. \textbf{204} (2026), 455--533. \href{https://doi.org/10.4007/annals.2026.204.2.1}{Article}.

\bibitem{ab} B. Adamczewski and Y. Bugeaud, \emph{Palindromic continued fractions}, Ann. Inst. Fourier \textbf{57} (2007), 1557--1574. \href{https://doi.org/10.5802/aif.2306}{doi:10.5802/aif.2306}.

\bibitem{ess-odd} C. Elsner, S. Shimomura, and I. Shiokawa, \emph{Algebraic relations for reciprocal sums of odd terms in Fibonacci numbers}, Ramanujan J. \textbf{17} (2008), 429--446. \href{https://doi.org/10.1007/s11139-007-9019-7}{Article}.

\bibitem{dnns-theta} D. Duverney, Ke. Nishioka, Ku. Nishioka, and I. Shiokawa, \emph{Transcendence of Jacobi's theta series}, Proc. Japan Acad. Ser.~A Math. Sci. \textbf{72} (1996), 202--203. \href{https://doi.org/10.3792/pjaa.72.202}{Article}.

\bibitem{bv} J. Bourgain and P. P. Varj\'u, \emph{Expansion in \(\SL_d(\Z/q\Z)\), \(q\) arbitrary}, Invent. Math. \textbf{188} (2012), 151--173. \href{https://doi.org/10.1007/s00222-011-0345-4}{doi:10.1007/s00222-011-0345-4}.

\bibitem{aperybb} J. M. Borwein and D. M. Bradley, \emph{Searching symbolically for Ap\'ery-like formulae for values of the Riemann zeta function}, ACM SIGSAM Bull. \textbf{30}(2) (1996), 2--7. \href{https://arxiv.org/abs/math/0505093}{Author version}.

\bibitem{ef} J.-H. Evertse and R. G. Ferretti, \emph{A further improvement of the Quantitative Subspace Theorem}, Ann. of Math. (2) \textbf{177} (2013), 513--590. \href{https://doi.org/10.4007/annals.2013.177.2.4}{doi:10.4007/annals.2013.177.2.4}.

\bibitem{serre} J.-P. Serre, \emph{Divisibilité de certaines fonctions arithmétiques}, Séminaire Delange--Pisot--Poitou \textbf{16} (1974/75). \href{https://www.numdam.org/item/SDPP_1974-1975__16_1_A15_0/}{Article}.

\bibitem{pellroth} K. F. Roth, \emph{Rational approximations to algebraic numbers}, Mathematika \textbf{2} (1955), 1--20. \href{https://doi.org/10.1112/S0025579300000644}{Article}.

\bibitem{mordell} L. J. Mordell, \emph{On Mr. Ramanujan's empirical expansions of modular functions}, Proc. Cambridge Philos. Soc. \textbf{19} (1917), 117--124.

\bibitem{tang} M. Tang, \emph{Relatively prime sets and a phi function for subsets of \(\{1,2,\ldots,n\}\)}, J. Integer Seq. \textbf{13} (2010), Art. 10.7.6. \href{https://cs.uwaterloo.ca/journals/JIS/VOL13/Tang/tang2.pdf}{Article}.

\bibitem{wolf} M. Wolf, \emph{Continued fractions constructed from prime numbers} (2010). \href{https://arxiv.org/abs/1003.4015}{arXiv:1003.4015}.

\bibitem{cor} P. Corvaja, \emph{An explicit version of the theorem of Roth--Ridout}, Rend. Sem. Mat. Univ. Politec. Torino \textbf{53} (1995), 253--260. \href{https://seminariomatematico.polito.it/rendiconti/cartaceo/53-3/253.pdf}{Article}.

\bibitem{cm} T. Chatterjee and S. Mandal, \emph{On transcendence of non-periodic continued fractions associated with modular forms and arithmetic functions}, Bol. Soc. Mat. Mex. \textbf{32} (2026), Art. 141. \href{https://doi.org/10.1007/s40590-026-00974-2}{Journal}. \href{https://arxiv.org/abs/2602.13300}{author version}.

\bibitem{bug2007} Y. Bugeaud, \emph{Extensions of the Cugiani--Mahler theorem}, Ann. Sc. Norm. Super. Pisa Cl. Sci. (5) \textbf{6} (2007), 477--498. \href{https://www.numdam.org/item/ASNSP_2007_5_6_3_477_0/}{Numdam}.

\bibitem{bug2013} Y. Bugeaud, \emph{Automatic continued fractions are transcendental or quadratic}, Ann. Sci. Éc. Norm. Supér. \textbf{46} (2013), 1005--1022. \href{https://doi.org/10.24033/asens.2208}{Article}.

\end{thebibliography}

\end{document}
